% Hoeffding 1963: Theorem 3, its lemmas/comparisons, and Sections 5--6.
\chapter{Variance and Dependent-Sampling Concentration}
\label{sec:p15-variance-dependent-concentration}

The width of an interval and the variance of a variable contain different
information. An upper-tail bound can use a finite variance and an upper
bound without requiring any lower bound. Dependence requires an additional
argument: multiplying marginal exponential moments is justified by joint
independence, not by vanishing covariances.
This chapter reconstructs Hoeffding's variance theorem and its proof,
the variance comparisons, and the complete dependent-sampling constructions
in Sections 5 and 6 \cite[Thm.~3, Secs.~3--6]{Hoeffding1963BoundedSums}.
The common-range results used below are in
\cref{sec:p15-bounded-concentration-learning}.

\section{An upper bound and a finite variance}
\label{sec:p15-upper-variance}

\begin{hbdefinition}{def:p15-upper-variance-setting}{The centered upper-bounded setting}
Let \(n\geq1\) be an integer and let \(X_1,\ldots,X_n\) be jointly
independent real random variables with
\(EX_i=0\), \(v_i=EX_i^2<\infty\), and \(X_i\leq b\) almost surely,
where \(b>0\). Write
\[
 \overline X=n^{-1}\sum_{i=1}^nX_i,\qquad
 \sigma^2=n^{-1}\sum_{i=1}^nv_i.
\]
Thus \(\operatorname{Var}(\overline X)=\sigma^2/n\).
No common distribution or common individual variance is assumed.
For noncentered variables the condition is
\(X_i-EX_i\leq b\); centering must be performed separately for each variable.
\end{hbdefinition}

\begin{hblemma}{thm:p15-upper-moment-extremal}{proved=annex:p15.upper-moment-extremal}{The exact one-variable moment bound}
If \(b>0\), \(EX=0\), \(EX^2=v<\infty\), \(X\leq b\), and \(h>0\), then
\[
 Ee^{hX}\leq M_{b,v}(h):=
 \frac{b^2}{b^2+v}e^{-hv/b}+
 \frac{v}{b^2+v}e^{hb}.
\]
For \(v>0\), equality holds for the two-point distribution
\[
 \Pr\{X=-v/b\}=\frac{b^2}{b^2+v},\qquad
 \Pr\{X=b\}=\frac{v}{b^2+v}.
\]
For \(v=0\), \(X=0\) almost surely and \(M_{b,0}=1\).
This is Hoeffding's Lemma 2; the annex supplies a complete
Hermite-quadratic proof of the moment bound.
\end{hblemma}

\begin{hblemma}{thm:p15-variance-moment-concavity}{proved=annex:p15.variance-moment-concavity}{Averaging unequal variances}
For fixed \(c>0\), define
\[
 f_c(u)=\log\frac{e^{-cu}+ue^c}{1+u},\qquad u\geq0.
\]
Then \(f_c''(u)<0\) and \(f_c'(u)>0\), including the right derivatives
at zero. Also \(f_c(0)=0\) and \(f_c'(0)=e^c-1-c\).
Consequently, in \cref{def:p15-upper-variance-setting}, for every \(h>0\),
\[
 E e^{h\sum_iX_i}
 \leq M_{b,\sigma^2}(h)^n
 \leq \exp\!\left\{\frac{n\sigma^2}{b^2}(e^{hb}-1-hb)\right\}.
\]
The first comparison is the concavity argument of Hoeffding's Lemma 3;
the second is its tangent-line consequence.
\end{hblemma}

\begin{hbtheorem}{thm:p15-variance-tail-chain}{proved=annex:p15.variance-tail-chain}{Hoeffding's variance bound and its simpler consequences}
Assume \cref{def:p15-upper-variance-setting} and \(\sigma^2>0\).
For \(0<t<b\), put
\[
\begin{split}
 B_1(n,t;b,\sigma^2)
 &=
 \left\{
 (1+bt/\sigma^2)^{-(\sigma^2+bt)/(b^2+\sigma^2)}
 (1-t/b)^{-b(b-t)/(b^2+\sigma^2)}
 \right\}^{n},\\
 \lambda&=bt/\sigma^2,\qquad \tau=nt/b,\\
 h_1(\lambda)&=(1+\lambda^{-1})\log(1+\lambda)-1,\\
 h_2(\lambda)&=\frac{\lambda}{2(1+\lambda/3)},\qquad
 h_3(\lambda)=\frac12\operatorname{arsinh}(\lambda/2).
\end{split}
\]
Here \(\operatorname{arsinh}z=\log(z+\sqrt{1+z^2})\).
The source's equations (2.8)--(2.14) give
\[
 \Pr\{\overline X\geq t\}\leq B_1
 \leq B_2:=e^{-\tau h_1(\lambda)}
 \leq e^{-\tau h_2(\lambda)},\qquad
 B_2\leq e^{-\tau h_3(\lambda)}.
\]
In particular the Bernstein-form consequence is
\[
 \Pr\{\overline X\geq t\}
 \leq \exp\!\left\{-\frac{nt^2}{2(\sigma^2+bt/3)}\right\}.
\]
The Bennett-form bound \(B_2\) and its two simpler consequences remain
valid for every \(t>0\).

At \(t=b\), the continuous extension of \(B_1\) is
\(\{\sigma^2/(b^2+\sigma^2)\}^n\); for \(t>b\) the probability is zero.
If \(\sigma^2=0\), every variable is zero almost surely and the probability
is zero for every \(t>0\). If the upper bound is \(b=0\), centering again
forces this degenerate case. An upper bound \(b<0\) is incompatible with
an integrable zero-mean variable.
For a two-sided deviation, apply the result also to \(-X_i\), which
requires its own upper bound, and add the two probabilities.
\end{hbtheorem}

\begin{hbcorollary}{thm:p15-variance-maximal}{proved=annex:p15.variance-maximal}{The source's independent partial-sum maximum}
Under \cref{def:p15-upper-variance-setting}, put
\(S_0=0\) and \(S_j=\sum_{i=1}^jX_i\). For \(h>0,t>0\),
\[
 \Pr\{\max_{0\leq j\leq n}S_j\geq nt\}
 \leq e^{-hnt}Ee^{hS_n}.
\]
All bounds and endpoint qualifications in
\cref{thm:p15-variance-tail-chain} therefore remain valid with
\(\{\overline X\geq t\}\) replaced by this maximum event.
The variance parameter remains \(n^{-1}\sum_i EX_i^2\).
The independence premise is retained; this corollary does not replace
it by a statement about the unconditional variances of dependent
martingale increments.
\end{hbcorollary}

\begin{hbtheorem}{thm:p15-variance-bound-comparison}{proved=annex:p15.variance-bound-comparison}{Exact comparison and the role of known variance}
For \(0<t<b\) and \(\sigma^2>0\), define
\[
 v=\frac{bt}{bt+\sigma^2},\quad w=\frac tb,\quad
 \rho(x)=\frac{(1-x)\log(1-x)+x-x^2/2}{x^2}\quad(0<x<1).
\]
Then
\[
 \frac{B_1}{B_2}=e^{-n\phi(v,w)},\qquad
 \phi(v,w)=\frac{w\{\rho(v)+\rho(w)\}}{v^{-1}+w^{-1}-1},
 \qquad
 \rho(x)=\sum_{j=1}^{\infty}\frac{x^j}{(j+1)(j+2)}.
\]
In particular
\[
 \frac{w(v+w)}{6(v^{-1}+w^{-1}-1)}
 <\phi(v,w)<
 \frac{w(v+w)}{2(v^{-1}+w^{-1}-1)}
 <\frac{t^3}{2b\sigma^2}.
\]
Thus \(e^{-nt^3/(2b\sigma^2)}<B_1/B_2<1\).
For fixed positive \(b,\sigma^2\), as \(t\downarrow0\),
\(\phi(v,w)=t^3/(6b\sigma^2)+O(t^4)\).
Both \(h_1(\lambda)\) and \(h_2(\lambda)\) equal
\(\lambda/2-\lambda^2/6+O(\lambda^3)\) as \(\lambda\downarrow0\).

The bound \(B_1\) is nondecreasing in \(\sigma^2\).
If also \(a\leq X_i\leq b\), where \(a<0<b\), then
\(\sigma^2\leq-ab\). Replacing \(\sigma^2\) by \(-ab\) in \(B_1\)
gives exactly the common-range exponential bound from
\cref{thm:p15-conc-common-range}, with normalized mean
\(-a/(b-a)\) and normalized displacement \(t/(b-a)\).
This bridge assumes the separately centered means of the present theorem;
it does not add a common-mean requirement to the source's Theorem 1.
\end{hbtheorem}

\begin{hbtheorem}{thm:p15-variance-optimality-boundary}{proved=annex:p15.variance-optimality-boundary}{What exponential optimality does and does not mean}
For fixed integer \(n\geq1\), \(b,\sigma^2>0\), and \(0<t<b\),
\(B_1\) is the minimum exponential-Markov
bound for the independent equal-variance two-point law in
\cref{thm:p15-upper-moment-extremal}. That law also attains the universal
moment bound at every positive \(h\). At \(t=b\) its tail probability
equals \(B_1\).

This does not make \(B_1\) the least possible uniform bound on the actual
tail for \(0<t<b\). For \(n=1,b=1,\sigma^2=1,t=1/2\),
\[
 B_1=\frac{2}{3^{3/4}}>\frac45.
\]
The sharp uniform tail bound at those values is \(4/5\), attained by
the law with masses \(4/5,1/5\) at \(1/2,-2\).

More generally, let \(W\) be real, let \(J\geq1\) be finite, and choose
real \(c_j,\alpha_j\) with \(Ee^{\alpha_jW}<\infty\).
If \(Q(s)=\sum_{j=1}^Jc_je^{\alpha_js}\) satisfies
\(Q(W)\geq0\) almost surely and \(Q(W)\geq1\) on \(\{W\geq a\}\),
then
\[
 \Pr\{W\geq a\}\leq\sum_{j=1}^Jc_jEe^{\alpha_jW}.
\]
This is the source footnote's more general exponential-majorant route.
A negative coefficient requires a lower moment bound, or the exact
moment, when bounding its contribution from above.

Among reductions of the form \(e^{-\tau H(\lambda)}\) obtained uniformly
from \(B_1\) for all admissible positive \(b,\sigma^2,t\), the largest
possible exponent function is \(H=h_1\). This last assertion concerns
that reduction of an exponential bound, not exact tail optimization.
\end{hbtheorem}

\section{Convex averages of independent blocks}
\label{sec:p15-dependent-blocks}

\begin{hblemma}{thm:p15-concentration-convex-mixture}{proved=annex:p15.concentration-convex-mixture}{A pointwise convex average transfers moment bounds}
Let \(T=\sum_{j=1}^Np_jT_j\), with deterministic \(p_j\geq0\),
\(\sum_jp_j=1\), and integrable real \(T_j\) on one probability space.
The \(T_j\) may be dependent. For every \(h>0\),
\[
 E e^{h(T-ET)}
 \leq\sum_jp_j E e^{h(T_j-ET_j)},\qquad
 \Pr\{T-ET\geq t\}
 \leq e^{-ht}\sum_jp_j E e^{h(T_j-ET_j)}.
\]
Terms with zero weight are omitted. The comparisons allow infinite right
sides. A common finite moment
bound for all the centered \(T_j\) therefore transfers to \(T\).
This is an average of variables on the same sample, not an assumption
that a random mixture selects one independent component.
\end{hblemma}

\begin{hbdefinition}{def:p15-ordered-u-statistic}{Ordered tuples and block averages}
For integers \(n\geq r\geq1\), put
\((n)_r=n(n-1)\cdots(n-r+1)\), and set \((m)_0=1\) for every integer
\(m\geq0\). Let \(X_1,\ldots,X_n\) be independent
random elements in a common measurable space \(\mathcal X\), and let
\(g:\mathcal X^r\to\mathbb R\) be a measurable function of \(r\)
arguments, integrable for every distinct-index tuple. Define
\[
 U=\frac1{(n)_r}\sum_{i_1,\ldots,i_r\ {\rm distinct}}
                   g(X_{i_1},\ldots,X_{i_r}),\qquad k=\lfloor n/r\rfloor.
\]
No symmetry of \(g\) is required: the average uses ordered tuples.
For a permutation \(\pi\) of \(\{1,\ldots,n\}\), put
\[
 V_\pi=\frac1k\sum_{j=1}^k
 g(X_{\pi((j-1)r+1)},\ldots,X_{\pi(jr)}).
\]
Each \(V_\pi\) averages \(k\) mutually independent block variables.
\end{hbdefinition}

\begin{hbtheorem}{thm:p15-one-sample-u-concentration}{proved=annex:p15.one-sample-u-concentration}{One-sample permutation representation}
In \cref{def:p15-ordered-u-statistic},
\[
 U=\frac1{n!}\sum_\pi V_\pi.
\]
If \(a\leq g\leq b\), \(a<b\), then for \(t>0\)
\[
 \Pr\{U-EU\geq t\}\leq e^{-2kt^2/(b-a)^2}.
\]
If the \(X_i\) are identically distributed, set
\(\mu=Eg(X_1,\ldots,X_r)=EU\).
For \(0\leq g\leq1\), all three common-range bounds and their endpoint
qualifications in \cref{thm:p15-conc-common-range} transfer with \(n=k\).
If \(g-\mu\leq b_0\), \(b_0>0\), and
\(\sigma_g^2=\operatorname{Var}(g(X_1,\ldots,X_r))<\infty\), the bounds
in \cref{thm:p15-variance-tail-chain} transfer with
\(n=k,b=b_0,\sigma^2=\sigma_g^2\).
The variance here is that of one kernel value, not of \(U\).
\end{hbtheorem}

\begin{hbtheorem}{thm:p15-two-sample-u-concentration}{proved=annex:p15.two-sample-u-concentration}{Two-sample permutation representation}
Let all \(X_1,\ldots,X_m,Y_1,\ldots,Y_n\) be mutually independent,
\(m\geq r\geq1,n\geq s\geq1\), and let \(g\) be a measurable integrable
kernel of \(r\) first-sample and \(s\) second-sample arguments.
All \(X_i\) take values in one measurable space and all \(Y_j\) in
another; the kernel is defined on the corresponding product. Define
\[
 U=\frac1{(m)_r(n)_s}
 \sum_{\substack{i_1,\ldots,i_r\ {\rm distinct}\\
                 j_1,\ldots,j_s\ {\rm distinct}}}
 g(X_{i_1},\ldots,X_{i_r},Y_{j_1},\ldots,Y_{j_s}),\quad
 k=\min(\lfloor m/r\rfloor,\lfloor n/s\rfloor).
\]
For permutations \(\pi,\eta\), let \(V_{\pi,\eta}\) average the \(k\)
kernel values formed from the successive disjoint \(r\)-blocks in
\(\pi\) and successive disjoint \(s\)-blocks in \(\eta\). Then
\[
 U=\frac1{m!n!}\sum_{\pi,\eta}V_{\pi,\eta}.
\]
If \(a\leq g\leq b\), \(a<b\), then for every \(t>0\) the upper-tail
bound for \(U-EU\) is \(e^{-2kt^2/(b-a)^2}\).
If each sample is identically distributed within itself, the common-range
and variance-sensitive conclusions of the preceding theorem hold with
this \(k\) and the mean and variance of one mixed kernel.
The two sample distributions need not be equal.
\end{hbtheorem}

\begin{hbtheorem}{thm:p15-repeated-index-concentration}{proved=annex:p15.repeated-index-concentration}{Repeated indices reduce to an explicit kernel average}
Let \(n\geq r\geq1\), with independent \(X_i\) in a common measurable
space and a real measurable kernel \(g\) on its \(r\)-fold product.
Define
\[
 W=\frac1{n^r}\sum_{i_1,\ldots,i_r=1}^n
                       g(X_{i_1},\ldots,X_{i_r}).
\]
Assume all displayed terms integrable, including repeated arguments.
For a partition \(\mathcal P\) of \(\{1,\ldots,r\}\) into \(j\) nonempty
blocks, fix a labeling of its blocks and let \(\beta_{\mathcal P}(\ell)\)
be the label of the block containing \(\ell\). Set
\[
 g^*(x_1,\ldots,x_r)=
 \sum_{\mathcal P}\frac{(n)_{|\mathcal P|}}{n^r}
 g(x_{\beta_{\mathcal P}(1)},\ldots,x_{\beta_{\mathcal P}(r)}).
\]
The coefficients sum to one and
\(W=U_r(g^*)\), the ordered \(U\) statistic with kernel \(g^*\).
Thus \(a\leq g\leq b\), \(a<b\), on all tuples implies, for \(t>0\),
\(\Pr\{W-EW\geq t\}\leq e^{-2\lfloor n/r\rfloor t^2/(b-a)^2}\).
For \(r=2\) one may use
\[
 g^*(x_1,x_2)=\frac{n-1}{n}g(x_1,x_2)+\frac1n g(x_1,x_1).
\]
For \(r=3\) one valid choice is
\[
\begin{split}
g^*={}&\frac{(n-1)(n-2)}{n^2}g(x_1,x_2,x_3)\\
&+\frac{n-1}{n^2}\{g(x_1,x_1,x_2)+g(x_1,x_2,x_1)+g(x_2,x_1,x_1)\}
+\frac1{n^2}g(x_1,x_1,x_1).
\end{split}
\]
These retain the source's equations (5.16)--(5.17).
\end{hbtheorem}

\begin{hbtheorem}{thm:p15-cramer-von-mises-kernel}{proved=annex:p15.cramer-von-mises-kernel}{The source's repeated-index goodness-of-fit example}
For an integer \(n\geq1\), a continuous cumulative distribution function
\(G\) on \(\mathbb R\)
and empirical distribution \(F_n(z)=n^{-1}\sum_i{\bf1}_{\{X_i\leq z\}}\),
define
\[
 \omega^2=\int_{\mathbb R}(F_n(z)-G(z))^2\,dG(z).
\]
For every observed sample,
\(\omega^2=n^{-2}\sum_{i,j=1}^n g(X_i,X_j)\), where
\[
 g(x,y)=\frac13+\frac12G(x)^2+\frac12G(y)^2-\max\{G(x),G(y)\}.
\]
For \(n\geq2\) this is the preceding repeated-index statistic of order two.
Its kernel range lies in \([-1/6,1/3]\). For independent \(X_i\),
\(n\geq2\), and \(t>0\),
\(\Pr\{\omega^2-E\omega^2\geq t\}\leq e^{-8\lfloor n/2\rfloor t^2}\).
This bound does not require \(X_i\) to have distribution \(G\).
For two observations with transformed values \(1/4,3/4\),
the statistic is \(1/48\).
\end{hbtheorem}

\begin{hbdefinition}{def:p15-finite-range-dependence}{Dependence between separated blocks}
A finite sequence of real random variables \(Y_1,\ldots,Y_n\)
is \((r-1)\)-dependent if
the vectors \((Y_1,\ldots,Y_i)\) and \((Y_j,\ldots,Y_n)\) are independent
whenever \(j-i\geq r\). Take \(1\leq r\leq n\).
If a stated dependence range has \(r>n\), replace it by \(r=n\);
both conditions then impose no restriction.
Let
\[
 I_i=\{i,i+r,i+2r,\ldots\}\cap\{1,\ldots,n\},\quad
 n_i=|I_i|=\left\lfloor\frac{n-i+r}{r}\right\rfloor,\quad
 S_i=\sum_{j\in I_i}Y_j,\quad k=\lfloor n/r\rfloor .
\]
\end{hbdefinition}

\begin{hbtheorem}{thm:p15-finite-range-concentration}{proved=annex:p15.finite-range-concentration}{Residue-class decomposition with unequal block sizes}
Assume the preceding dependence condition and integrable \(Y_i\).
Within each \(I_i\) the variables are jointly independent, and, for
\(S=\sum_jY_j\),
\[
 \frac{S-ES}{n}
 =\sum_{i=1}^r\frac{n_i}{n}\frac{S_i-ES_i}{n_i}.
\]
For every \(h>0\) and \(t>0\),
\[
 \Pr\{S-ES\geq nt\}
 \leq\sum_{i=1}^r\frac{n_i}{n}e^{-ht}
       E\exp\{h(S_i-ES_i)/n_i\}.
\]
If \(a\leq Y_i\leq b\), \(a<b\), the bound is
\(e^{-2kt^2/(b-a)^2}\). If in addition the \(Y_i\) have a common
distribution, the common-range bounds for \(0\leq Y_i\leq1\), or
the variance bounds for \(Y_i-EY_i\leq b_0\) with finite variance,
transfer with effective count \(k\).
This remains true when the positive sizes \(n_i\) are unequal.
\end{hbtheorem}

\section{Sampling without replacement}
\label{sec:p15-finite-population-concentration}

\begin{hbdefinition}{def:p15-finite-population-sampling}{The two finite-population sampling laws}
Fix an integer \(N\geq1\) and real values \(c_1,\ldots,c_N\), allowing
repetitions of values.
A sample \(X_1,\ldots,X_n\), \(1\leq n\leq N\), without replacement
uses a uniformly chosen ordered tuple of distinct population indices.
A sample \(Y_1,\ldots,Y_n\) with replacement uses independent uniform
indices from \(\{1,\ldots,N\}\). Write
\[
 \mu=\frac1N\sum_jc_j,\qquad
 \sigma_C^2=\frac1N\sum_j(c_j-\mu)^2.
\]
Repetition of a population value is different from repetition of its index.
\end{hbdefinition}

\begin{hbtheorem}{thm:p15-without-replacement-convex}{proved=annex:p15.without-replacement-convex}{Hoeffding's finite-population convex comparison}
For the two sampling laws and every real continuous convex \(f\),
\[
 E f\!\left(\sum_{i=1}^nX_i\right)
 \leq E f\!\left(\sum_{i=1}^nY_i\right).
\]
Consequently the common-range bounds of
\cref{thm:p15-conc-common-range}, the common-interval sum specialization of
\cref{thm:p15-conc-unequal-ranges}, and the variance bounds of
\cref{thm:p15-variance-tail-chain} for an independent sample from this
population also bound the corresponding unweighted without-replacement
sample mean.
Use the population interval, mean \(\mu\), and variance \(\sigma_C^2\)
in those bounds. In the variance theorem use centered population values
\(c_j-\mu\) and an upper bound for those values.

Both sample means have expectation \(\mu\). For \(N>1\),
\[
 \operatorname{Var}(\overline X)
 =\frac{N-n}{N-1}\frac{\sigma_C^2}{n},\qquad
 \operatorname{Var}(\overline Y)=\frac{\sigma_C^2}{n}.
\]
For \(N=1\) both variances are zero. The proved exponential comparison
does not replace its parameter \(\sigma_C^2\) by the smaller
finite-population-corrected variance. At \(n=N\) the first sample mean is
deterministic even though the transferred upper bound can be positive.
\end{hbtheorem}

\begin{hbtheorem}{thm:p15-dependent-concentration-examples}{proved=annex:p15.dependent-concentration-examples}{Worked blocks, finite samples and a pairwise-independence failure}
The following calculations illustrate the preceding results. Tail bounds
in these examples use \(t>0\).
\begin{enumerate}
\item Let twenty independent variables take values \(1,-1/4\) with
probabilities \(1/5,4/5\). Their means are zero and variances \(1/4\).
At \(t=1/4\), the exact upper tail and four bounds are
\[
\begin{split}
 \Pr\{\overline X\geq1/4\}
 &=5^{-20}\sum_{j=8}^{20}\binom{20}{j}4^{20-j},\\
 B_1&=2^{16}/3^{12},\qquad B_2=e^5/2^{10},\\
 B_{\rm Bernstein}&=e^{-15/8},\qquad
 B_{\rm Prohorov}=e^{-(5/2)\operatorname{arsinh}(1/2)} .
\end{split}
\]
The exact tail is approximately \(0.03214\); the bounds are approximately
\(0.12332,0.14493,0.15335,0.30028\), respectively.
\item Let \(\Pr\{Z=2^j\}=7\,8^{-j}\), \(j=1,2,\ldots\), and
\(X=7/3-Z\). Then \(EX=0\), \(EX^2=14/9\), and \(X\leq1/3\).
This law meets the one-variable upper-tail premises and has no lower
bound. For every \(h>0\), \(Ee^{-hX}=\infty\); the available upper-tail
moment estimate cannot be reused as a lower-tail estimate.
\item For independent identically distributed observations
\(Z_i=(A_i,B_i)\), the order-two kernel
\(g(Z_i,Z_j)=\operatorname{sign}((A_i-A_j)(B_i-B_j))\), with
\(\operatorname{sign}(0)=0\), has range \([-1,1]\).
Its ordered average is the unadjusted Kendall pair statistic and has upper-tail bound
\(e^{-\lfloor n/2\rfloor t^2/2}\). For the three observations
\((1,1),(2,3),(3,2)\), its value is \(1/3\).
\item For independent samples, the proportion of ordered pairs satisfying
\(Y_j<X_i\) uses \(g(x,y)={\bf1}_{\{y<x\}}\), with \(r=s=1\).
This is the normalized Wilcoxon--Mann--Whitney pair count.
Its upper-tail bound is \(e^{-2\min(m,n)t^2}\).
For \(X=(1,3)\), \(Y=(0,2,4)\), the observed proportion is \(3/6=1/2\).
\item If \(Z_1,\ldots,Z_{n+r-1}\) are independent and
\(Y_i=Z_i\cdots Z_{i+r-1}\), then \(Y_i\) is \((r-1)\)-dependent.
For independent Bernoulli\((p)\) inputs, \(0\leq Y_i\leq1\),
\(EY_i=p^r\), and its empirical mean has the common-range bounds with
effective count \(\lfloor n/r\rfloor\).
For \(n=5,r=2\), the two residue sets have sizes \(3,2\), with weights
\(3/5,2/5\), and the range-only bound is \(e^{-4t^2}\).
\item From population \((0,0,1,1)\), sample two indices.
The without-replacement sum has probabilities \((1/6,4/6,1/6)\)
at \(0,1,2\); the with-replacement sum has probabilities
\((1/4,1/2,1/4)\). Their mean variances are \(1/12\) and \(1/8\).
The with-replacement exponential moment minus the without-replacement
one is \((e^h-1)^2/12\). At displacement \(t=1/2\), the exact
without-replacement upper tail is \(1/6\); the transferred endpoint
bound is \(1/4\).
\item If \(A,B\) are independent fair random signs, the three variables
\(A,B,AB\) are pairwise independent, centered, upper bounded by one,
and each has variance one. Their mean equals one with probability
\(1/4\), exceeding the independent variance-theorem endpoint bound
\(2^{-3}=1/8\). Pairwise independence does not suffice.
\end{enumerate}
\end{hbtheorem}

The derived bounds are finite-sample probability statements under the
specified joint laws. They do not infer a sampling law from a table of
observations. In particular overlapping kernel terms and adjacent blocks
cannot be counted as mutually independent observations.
